\section{Implemented workflow and authority boundaries}
The application combines a Tauri/Rust backend, Monaco source editing and a
PDF.js preview. The Rust document model owns canonical source/history and
visual-operation authority. JavaScript provides interaction and prepares the
view but cannot manufacture a native approval. Figure~\ref{fig:workflow}
summarizes the supported path; the following account separates selection,
candidate verification, commit and durability.

\begin{figure}[htbp]
\centering\small
\newlength{\stagewidth}\setlength{\stagewidth}{.27\textwidth}
\begin{tabularx}{\linewidth}{@{}YcYcY@{}}
\fbox{\parbox{\stagewidth}{\centering\textbf{1. Select}\\Compiled text hit/group}} & $\to$ &
\fbox{\parbox{\stagewidth}{\centering\textbf{2. Map}\\Source hints and unit}} & $\to$ &
\fbox{\parbox{\stagewidth}{\centering\textbf{3. Establish identity}\\Private zero-offset probe}}\\[9pt]
\fbox{\parbox{\stagewidth}{\centering\textbf{4. Plan}\\Immutable source candidate}} & $\to$ &
\fbox{\parbox{\stagewidth}{\centering\textbf{5. Compile and inspect}\\Operation-specific policy}} & $\to$ &
\fbox{\parbox{\stagewidth}{\centering\textbf{6. Decide}\\Commit or refusal}}
\end{tabularx}
\caption{Read stages in numerical order. Selection capability, native commit,
view publication and later Save are distinct boundaries.}
\label{fig:workflow}
\end{figure}

\subsection{Visible selection and source evidence}
PDF.js supplies compiled text runs, line geometry and marked-content bounds.
The canvas controller proposes a clicked line or adjacent group. A marquee
does not directly choose source offsets or write files. Read-text mode remains
available separately for ordinary highlighting and copying.

A later source-only UI revision links an accepted PDF selection to its
normalized source range using a Monaco decoration, preserving canvas focus.
It also adds a distinct reset-to-zero action that retains the movable wrapper,
separate from removing positioning. Literal editing shows bounded formatting,
direction and protected-command hints; these are not a TeX font/shaping preview.
The modal is placed beside the selected unit only when it fits the viewport,
with a centered fallback. The original source-only receipt did not establish
Windows usability. Later builds and the separately recorded owner observations
have narrower, named scopes; source presence alone remains insufficient.

For unwrapped text, the resolver asks native SyncTeX for file/line evidence
and examines conservative source units. It can enumerate bounded alternatives;
the frontend limits distinct probe attempts rather than searching indefinitely.
Duplicate hints are not repeatedly compiled. A wrapped paragraph may be the
smallest safe unit, even if the user clicked one displayed line. Cross-file,
incompatible-column and unsupported structural evidence can refuse before
any edit is planned. A repeated printed string is not treated as a unique
source occurrence merely because its spelling matches.

An existing movable object is resolved through its compiled marker and
observed bounds. Multiple PDF instances, multiple pages or unsupported page
coordinates do not grant an ordinary movement frame. The marker links a
source object to emitted content; it does not independently establish that
the user intended that object.

\subsection{Private identity probes}
For an ordinary source unit, a private candidate wraps the fragment at zero
displacement. The frontend requires compatible pages, unchanged runs and
pixels, unique marker membership and coverage of the clicked runs. The
selection can expand to the verified unit, and the frame shows that scope.
The probe does not add a wrapper to canonical source or history. Its cached
proof belongs to the current baseline view and must still pass baseline
dependency checks when reused.

A probe token is not a visual-operation approval. The eventual movement
requires a native operation plan and independent native layout approval.
First-wrap movement also performs a native zero-offset identity check. Thus
frontend selection evidence and native commit authority serve different
purposes. Their agreement is useful, but shared source assumptions and the
movement package can still create correlated mistakes.

\subsection{Exact source planning}
Canonical sources preserve UTF-8 bytes, including BOMs and local line endings.
Monaco sees an LF-normalized view without a leading BOM and reports UTF-16
coordinates. The Rust projection maps valid editor boundaries back to source
byte positions. It does not serialize a syntax tree over the whole document.
Unsupported or stale ranges can refuse rather than guess through a surrogate
or source boundary.

For supported text, native planning derives the candidate from the current
sources, object/selection, expected content and operation. A proposal with an
unrequested neighboring or supporting-file change does not match this plan.
The movement macro applies drawing translation while retaining the fragment's
layout contribution. It does not generally box and relocate arbitrary material
outside its formatting context or solve a new flow layout. Support-package
version consistency is also checked.

The frontend's proposed source and typed messages are not the final authority.
Rust-generated protocol contracts and runtime reply validation protect the
transport boundary; native prepare still rederives the permitted source
change. Reusing the same planner at two boundaries is not equivalent to two
independent semantic specifications, which motivates a separate source oracle
in the evaluation.

\subsection{Immutable builds and native inspection}
Candidates compile from private snapshots, with engine, recipe, source,
asset and recorded dependency receipts. The current direct-engine service
has bounded passes/time and declared bibliography recipes. Preflight failure,
truncated output or unresolved build diagnostics can withhold verification
eligibility. Source remains editable when TeX compilation fails; the last
good PDF is retained as an older view rather than promoted to current authority.
Shell escape is disabled and MiKTeX automatic package installation is disabled.
These process controls are not a security sandbox for arbitrary untrusted TeX.

For a visual candidate with a native operation plan, the build RPC obtains independent inspection
from a checked, pinned PDFium worker before returning successful visual
approval. Normal displacement compares against a baseline-derived translation
prediction, including target geometry, unaffected material and complete
same-renderer appearance. Candidate pixels do not construct that prediction.
Budget, paint, annotation, target-isolation or correspondence failures can refuse.

After the native reply, the frontend prepares the PDF offscreen in PDF.js,
checks its operation-specific geometry and rechecks the build receipt.
The latter native \code{verify\_build} call checks recorded inputs; it is not
another independent PDF comparison. This ordering matters when interpreting
timings: native layout approval is already included in the build RPC.

\subsection{Commit, cancellation and newer typing}
Approval is private native state bound to the exact project, revision, build
and cancellation lease. Commit checks it under the canonical document lock
and consumes the pending candidate once. A JSON field cannot substitute for
that approval. Source edits, changed configuration, stale results and cancelled
leases can invalidate an operation even when a previously compiled PDF exists.

The frontend also checks identity after asynchronous steps. If the user types
while a visual candidate compiles, the newer source is preserved and the old
candidate is discarded. Preparation failure destroys the private prepared
view rather than adding history. Dragging updates an outline and does not
compile on pointer movement; release requests an edit. Event-boundary tests
cover Escape, capture loss and focus-loss cancellation. The owner later reported
the named r11 held-drag cancellation checklist passed. That guided observation
does not cover every capture-loss interleaving or Windows display transition.

Native commit occurs before final frontend publication. Offscreen preparation
reduces the chance of a presentation failure after commit, but this cross-process
workflow is not claimed as a filesystem/UI-wide atomic transaction. Native
acceptance and a displayed frame are recorded separately in the proposed study.

The subsequent failure-handling revision distinguishes a known
native commit followed by presentation/reselection failure from a refusal
before commit. If an attempted commit has no usable acknowledgement or
confirming state, the replica stops new writes and records uncertainty rather
than automatically retrying. Cleanup failures must not clear the busy state
incorrectly or destroy an adopted PDF. Resumed controller tests exercise
acknowledged-commit publication failure, unknown acknowledgement and control
release; they are not injected failures in a running Windows process. Native
history, view publication, original-file Save and exact reopen still require
separate observations.

\subsection{Source-model extension under development}
The 10 November core adds typed source-byte and editor-UTF-16 coordinates,
logical insertion affinity and transient change maps. Unmodified selections
map through preceding changes; replacement through a selected interval clears
identity instead of retaining numerically valid offsets over different text.
All secondary ranges must survive, and recorded history restores the selection
with its source snapshot. Scalar boundaries are not grapheme or visual caret
stops. Existing entity references remain session/revision-bound and become
stale after any changed revision; the coordinate map does not revive them or
revalidate old rendered geometry.

A conservative source projection retains unsupported syntax as opaque bytes.
The first typed intent planner prepares literal paragraph replacement and
paragraph-content deletion against exact current references. It emits ordinary
\LaTeX, changes one rederived source range and preserves other files/profile.
Escaped input is distinct from raw source; structural input requires a different
intent. The plan records a rendered-content proof obligation and declared
document reflow, with no visual-acceptance method. Native rendered verification,
source-to-PDF identity and UI integration remain open. The later native draft
transport has a separately bounded acknowledgement contract described below.
The projection can decline an entire macro-defined document, including the
current template, so useful editable coverage is unresolved. This first slice
is not the full Word-like authoring workflow.

The later core at \code{897c121} recognizes narrow preamble-only forms for
plain font names, the current default languages, specific Arabic/Hebrew font
switches and literal paragraph dimensions. Unknown options, computed arguments,
general macro definitions and body context changes still refuse. Recognition
preserves source bytes; it does not establish package semantics or rendered
identity. Current syntax references are the maintained fontspec user-command
documentation and Polyglossia manual~\cite{fontspeccommands,polyglossiamanual}.

The explicit \code{apply\_source\_intent\_draft} core API consumes a privately
constructed plan through existing source/history authority. Its receipt records
the source revision and a still-required render proof. Foreign/stale plans and
revision exhaustion cannot mutate source; no-op drafts retain selection/history.
This is the source-editor lifecycle, requiring journaling and invalidation of
rendered editing capabilities at the integration layer. The core API alone
provides no IPC/UI handler, visual approval or automatic Save. The r10
navigation artifact predates these later core additions.

\paragraph{Native draft transport and acknowledgement resolution.}
The 13 November continuation supplies native-issued, session/revision-bound
paragraph handles and strict tagged Replace/Delete draft requests. Caller
offsets cannot substitute for those private references. The native service
rederives the core plan, validates its proposed history state and protects
changed drafts before publishing canonical memory. Pending visual candidates
and obsolete compilation are invalidated only for actual changes. This adds
neither a PDF-selection command nor rendered approval.

An operation ID identifies one exact request within an open session. Up to 128
receipts and 4 MB of request payload are retained without eviction; capacity
refuses new requests rather than forgetting a destructive outcome. A duplicate
returns the immutable original outcome and a separately current snapshot, so
retry after Undo cannot silently redo the operation. A conflicting payload
refuses. Status lookup performs no source edit.

Receipts and handles do not survive session reopen. Unknown lookup then means
that this session has no receipt, not that the previous process did not act.
The caller must reconcile current/recovered source instead of automatically
issuing a new request. A journal-write error can also follow partial filesystem
writes: canonical memory remains unchanged, but the cached outcome is explicitly
uncertain rather than claiming no draft could later recover. Thus the contract
is bounded session idempotency and protected draft recovery, not durable
exactly-once execution across crashes. No automatic original-file Save follows.

\paragraph{Admitted refusal is a definite source outcome.}
The 14 November extension returns and caches \code{source\_draft\_refused}
for an admitted bounded request rejected before canonical or recovery mutation.
It binds the requested and observed revisions and a finite reason. A stale
revision, unavailable target or unsupported literal can therefore resolve
without reopening the project. An identical retry retains the original refusal
even after later source changes; lookup still returns separately current state.
Selection, redo, pending candidates and running work are preserved by refusal.

This contract does not convert every exception into nonexecution. Schema/ID,
missing-session, conflicting-payload and capacity failures remain outside the
admitted receipt contract. An unknown lookup cannot clear an unresolved call,
and a recovery-write failure stays uncertain. Table~\ref{tab:source-outcomes}
states the permitted response; neither known outcome grants rendered approval.

\begin{table}[htbp]
\small\begin{tabularx}{\linewidth}{@{}>{\raggedright\arraybackslash}p{.19\linewidth}YY@{}}
\toprule
Confirmed outcome & What it establishes & Permitted client response\\
\midrule
Applied draft & Named source/history step; render proof unresolved & Adopt current state; clear matching uncertainty without reapplying\\
Admitted refusal & No canonical/recovery mutation by that request & Retain current source; allow a later edit after confirmation\\
Unknown / lost reply & No confirmed execution outcome in this session & Stop dependent writes; lookup/reconcile; no automatic retry\\
Protection uncertain & Live source retained, but a proposed recovery draft may exist & Preserve uncertainty; explicit current/recovered-state reconciliation\\
\bottomrule
\end{tabularx}
\caption{Outcome knowledge is separate from source revision, presentation and
original-file persistence.}\label{tab:source-outcomes}
\end{table}

\subsection{History, Save and recovery}
Table~\ref{tab:authority-levels} separates source progress from the stronger
claims that a later stage would need. A new transport method changes how a
draft is requested; it does not itself establish rendered identity or Save.

\begin{table}[htbp]
\small\begin{tabularx}{\linewidth}{@{}>{\raggedright\arraybackslash}p{.19\linewidth}YYY@{}}
\toprule
Entry / outcome & Canonical source/history & Rendered authority & Persistence\\
\midrule
Read-only typed plan & Unchanged & Obligation only & None\\
Explicit source draft & Exact guarded shared-history change & Unresolved; last-good view remains older & Protected draft; original files unchanged\\
Transport receipt or lookup & Records an outcome; lookup performs no edit & No PDF target/approval evidence & Session receipt; unknown after reopen\\
Compiled candidate & Not a commit & Parse/compile success alone insufficient & Private build artifacts\\
Approved visual transaction & Consumes exact native candidate & Named policy and current dependencies & History/recovery; still separate from Save\\
Explicit Save & Current source retained & No additional render proof & Guarded original-file replacement\\
\bottomrule
\end{tabularx}
\caption{Authority is specific to an entry layer and outcome, rather than a
single success flag.}\label{tab:authority-levels}
\end{table}

Source edits and accepted visual edits share native Undo/Redo history. Typing
may coalesce; an accepted visual transaction has its own entry. Undo restores
source state while advancing revision and invalidating old layout authority.
It does not authorize a stale candidate merely because old bytes reappear.

Save is explicit. It verifies external file fingerprints and stages multi-file
replacement with recovery information. Recovery can finish a staged Save,
restore earlier files or preserve the editor version as a separate copy.
Rollback and removal of a staged recovery journal are different operations.
Typing recovery protects unsaved source/history separately and grants no old
visual approval. These are recoverable writes, not a claim of one atomic
filesystem rename covering an entire project.

The code anchors are \path{web/app.js}, the modules under \path{web/core/},
\path{src-tauri/src/service/document.rs}, \path{src-tauri/src/service/layout_gate.rs},
and the document/build crates. \path{research/implementation-audit.md} gives
the detailed mapping and separates implemented, recorded and planned behavior.
